\documentclass[a4paper,12pt]{article}

\usepackage[T2A]{fontenc}
\usepackage[utf8]{inputenc}
\usepackage[russian]{babel}

\usepackage{amsmath, amssymb, amsthm, mathtools}
\usepackage{helvet}
\renewcommand{\familydefault}{\sfdefault}
\usepackage{icomma}
\usepackage{geometry}
\usepackage{microtype}
\usepackage{tikz}
\usepackage{tkz-euclide}
\usepackage{pgfplots}
\pgfplotsset{compat=1.18}
\usepackage{tabularx, booktabs, longtable}
\usepackage{enumitem}
\usepackage{hyperref}
\usepackage{parskip}
\usepackage{fancyhdr}
\usepackage{setspace}
\usepackage{xcolor}

\geometry{left=18mm,right=18mm,top=18mm,bottom=20mm}
\onehalfspacing
\setlist[itemize]{leftmargin=1.4em, itemsep=0.2em, topsep=0.2em}
\setlist[enumerate]{leftmargin=1.6em, itemsep=0.35em, topsep=0.35em}

\definecolor{accent}{HTML}{008080}
\definecolor{darkaccent}{HTML}{004D4D}
\definecolor{textgray}{HTML}{333333}
\definecolor{softgray}{HTML}{F4F6F6}

\hypersetup{
    colorlinks=true,
    linkcolor=darkaccent,
    urlcolor=darkaccent
}

\pagestyle{fancy}
\fancyhf{}
\lhead{\textcolor{textgray}{Чек-лист}}
\rhead{\textcolor{textgray}{Тригонометрические уравнения}}
\cfoot{\thepage}
\renewcommand{\headrulewidth}{0.4pt}
\renewcommand{\headrule}{\hbox to\headwidth{\color{accent}\leaders\hrule height \headrulewidth\hfill}}

\newcommand{\topic}[1]{\textcolor{darkaccent}{\textbf{#1}}}
\newcommand{\important}[1]{\textcolor{darkaccent}{\textbf{#1}}}
\newcommand{\Z}{\mathbb{Z}}

\begin{document}

\begin{titlepage}
\thispagestyle{empty}
\centering

\vspace*{28mm}

{\Huge\bfseries Чек-лист\par}
\vspace{6mm}
{\Large\bfseries Первые тригонометрические уравнения\par}
\vspace{3mm}
{\large\bfseries Единичная окружность, базовые случаи и формулы преобразования\par}

\vspace{22mm}

{\large Лёвин Артём Александрович эксклюзивно для Тимофея Васильченко\par}

\vspace{10mm}

{\large \today\par}

\vfill

\begin{tikzpicture}[remember picture, overlay]
    \draw[accent, line width=1.2pt]
        ([xshift=25mm,yshift=25mm]current page.south west)
        .. controls ([xshift=65mm,yshift=42mm]current page.south west)
        and ([xshift=-65mm,yshift=8mm]current page.south east)
        .. ([xshift=-25mm,yshift=28mm]current page.south east);
    \draw[darkaccent, line width=0.5pt]
        ([xshift=35mm,yshift=18mm]current page.south west)
        .. controls ([xshift=78mm,yshift=32mm]current page.south west)
        and ([xshift=-78mm,yshift=18mm]current page.south east)
        .. ([xshift=-35mm,yshift=20mm]current page.south east);
\end{tikzpicture}

\vspace*{18mm}

\end{titlepage}

\section*{1. Главная идея}

Тригонометрические уравнения часто сводятся к простым уравнениям вида
\[
    \sin x = -1,\quad \sin x=0,\quad \sin x=1,
\]
\[
    \cos x=-1,\quad \cos x=0,\quad \cos x=1.
\]

\topic{Ключевой принцип:} сначала упростить выражение с помощью формул, затем решить базовое тригонометрическое уравнение на единичной окружности.

\section*{2. Единичная окружность}

На единичной окружности точка, соответствующая углу \(x\), имеет координаты
\[
    (\cos x;\sin x).
\]

Значит:
\[
    \cos x \text{ --- это } x\text{-координата точки},
\]
\[
    \sin x \text{ --- это } y\text{-координата точки}.
\]

\begin{center}
\begin{tikzpicture}[scale=1.15]
    \draw[->, textgray] (-1.45,0) -- (1.55,0) node[right] {\(x\)};
    \draw[->, textgray] (0,-1.45) -- (0,1.55) node[above] {\(y\)};
    \draw[accent, line width=1pt] (0,0) circle (1);
    \fill[darkaccent] (1,0) circle (1.7pt) node[below right] {\(0\)};
    \fill[darkaccent] (0,1) circle (1.7pt) node[above right] {\(\frac{\pi}{2}\)};
    \fill[darkaccent] (-1,0) circle (1.7pt) node[below left] {\(\pi\)};
    \fill[darkaccent] (0,-1) circle (1.7pt) node[below right] {\(-\frac{\pi}{2}\)};
    \draw[->, darkaccent, line width=0.8pt] (0.35,0) arc (0:65:0.35);
    \node[darkaccent] at (0.45,0.32) {\(x\)};
    \node[textgray] at (0,-1.85) {один полный оборот равен \(2\pi\)};
\end{tikzpicture}
\end{center}

\topic{Важно:} окружность можно обходить сколько угодно раз. Поэтому в ответах появляется добавка:
\[
    2\pi n,\qquad n\in\Z.
\]

\section*{3. Базовые уравнения для синуса}

\renewcommand{\arraystretch}{1.45}
\begin{table}[h!]
\caption{Частные случаи для синуса}
\centering
\begin{tabularx}{\linewidth}{@{}lX@{}}
\toprule
\textbf{Уравнение} & \textbf{Общее решение} \\
\midrule
\(\sin x=1\) &
\[
    x=\frac{\pi}{2}+2\pi n,\qquad n\in\Z
\]
\\
\(\sin x=-1\) &
\[
    x=-\frac{\pi}{2}+2\pi n,\qquad n\in\Z
\]
\\
\(\sin x=0\) &
\[
    x=\pi n,\qquad n\in\Z
\]
\\
\bottomrule
\end{tabularx}
\end{table}

\topic{Почему для \(\sin x=0\) можно писать \(x=\pi n\)?}

Потому что подходят две точки окружности:
\[
    x=0+2\pi n
\]
и
\[
    x=\pi+2\pi n.
\]
Обе серии объединяются в одну:
\[
    x=\pi n,\qquad n\in\Z.
\]

\section*{4. Базовые уравнения для косинуса}

\begin{table}[h!]
\caption{Частные случаи для косинуса}
\centering
\begin{tabularx}{\linewidth}{@{}lX@{}}
\toprule
\textbf{Уравнение} & \textbf{Общее решение} \\
\midrule
\(\cos x=1\) &
\[
    x=2\pi n,\qquad n\in\Z
\]
\\
\(\cos x=-1\) &
\[
    x=\pi+2\pi n,\qquad n\in\Z
\]
\\
\(\cos x=0\) &
\[
    x=\frac{\pi}{2}+\pi n,\qquad n\in\Z
\]
\\
\bottomrule
\end{tabularx}
\end{table}

\topic{Почему для \(\cos x=0\) можно писать \(x=\frac{\pi}{2}+\pi n\)?}

Подходят две точки:
\[
    x=\frac{\pi}{2}+2\pi n
\]
и
\[
    x=-\frac{\pi}{2}+2\pi n.
\]
Они объединяются в одну серию:
\[
    x=\frac{\pi}{2}+\pi n,\qquad n\in\Z.
\]

\section*{5. Формулы сложения и вычитания}

Формулы нужно знать наизусть: без них трудно увидеть способ преобразования.

\[
    \sin(\alpha+\beta)=\sin\alpha\cos\beta+\cos\alpha\sin\beta,
\]
\[
    \sin(\alpha-\beta)=\sin\alpha\cos\beta-\cos\alpha\sin\beta.
\]

\[
    \cos(\alpha+\beta)=\cos\alpha\cos\beta-\sin\alpha\sin\beta,
\]
\[
    \cos(\alpha-\beta)=\cos\alpha\cos\beta+\sin\alpha\sin\beta.
\]

\topic{Главная идея:} если в уравнении есть произведения вида
\[
    \cos\alpha\cos\beta+\sin\alpha\sin\beta,
\]
то нужно увидеть формулу:
\[
    \cos(\alpha-\beta).
\]

А если есть
\[
    \cos\alpha\cos\beta-\sin\alpha\sin\beta,
\]
то это
\[
    \cos(\alpha+\beta).
\]

\section*{6. Пример с формулой косинуса разности}

Решим уравнение:
\[
    1-\cos 3x\cos 2x=\sin 3x\sin 2x.
\]

Переносим всё влево:
\[
    1-\cos 3x\cos 2x-\sin 3x\sin 2x=0.
\]

Выносим минус перед суммой:
\[
    1-(\cos 3x\cos 2x+\sin 3x\sin 2x)=0.
\]

По формуле
\[
    \cos(\alpha-\beta)=\cos\alpha\cos\beta+\sin\alpha\sin\beta
\]
получаем:
\[
    1-\cos(3x-2x)=0.
\]

Значит,
\[
    1-\cos x=0,
\]
\[
    \cos x=1.
\]

Ответ:
\[
    \boxed{x=2\pi n,\qquad n\in\Z.}
\]

\section*{7. Пример с раскрытием \(\sin\left(\frac{\pi}{4}-x\right)\)}

Решим уравнение:
\[
    2\sin\left(\frac{\pi}{4}-x\right)+\sqrt2\sin x=-\sqrt2.
\]

Раскрываем синус разности:
\[
    \sin\left(\frac{\pi}{4}-x\right)
    =
    \sin\frac{\pi}{4}\cos x-\cos\frac{\pi}{4}\sin x.
\]

Так как
\[
    \sin\frac{\pi}{4}=\cos\frac{\pi}{4}=\frac{\sqrt2}{2},
\]
то
\[
    \sin\left(\frac{\pi}{4}-x\right)
    =
    \frac{\sqrt2}{2}\cos x-\frac{\sqrt2}{2}\sin x.
\]

Подставляем:
\[
    2\left(\frac{\sqrt2}{2}\cos x-\frac{\sqrt2}{2}\sin x\right)+\sqrt2\sin x=-\sqrt2.
\]

Получаем:
\[
    \sqrt2\cos x-\sqrt2\sin x+\sqrt2\sin x=-\sqrt2.
\]

Слагаемые с \(\sin x\) сокращаются:
\[
    \sqrt2\cos x=-\sqrt2.
\]

Делим на \(\sqrt2\):
\[
    \cos x=-1.
\]

Ответ:
\[
    \boxed{x=\pi+2\pi n,\qquad n\in\Z.}
\]

\section*{8. Формулы двойного угла}

\[
    \sin 2\alpha=2\sin\alpha\cos\alpha.
\]

\[
    \cos 2\alpha=\cos^2\alpha-\sin^2\alpha.
\]

\topic{Особенно важно:}
если есть произведение
\[
    2\sin\alpha\cos\alpha,
\]
его можно заменить на
\[
    \sin 2\alpha.
\]

Если есть разность квадратов:
\[
    \cos^2\alpha-\sin^2\alpha,
\]
её можно заменить на
\[
    \cos 2\alpha.
\]

\section*{9. Пример с разностью четвёртых степеней}

Решим уравнение:
\[
    \cos^4 5x-\sin^4 5x=0.
\]

Используем формулу разности квадратов:
\[
    a^2-b^2=(a-b)(a+b).
\]

Получаем:
\[
    \cos^4 5x-\sin^4 5x
    =
    (\cos^2 5x-\sin^2 5x)(\cos^2 5x+\sin^2 5x).
\]

По основному тригонометрическому тождеству:
\[
    \cos^2 5x+\sin^2 5x=1.
\]

Значит:
\[
    \cos^2 5x-\sin^2 5x=0.
\]

Используем формулу:
\[
    \cos 2\alpha=\cos^2\alpha-\sin^2\alpha.
\]

При \(\alpha=5x\):
\[
    \cos 10x=0.
\]

Тогда:
\[
    10x=\frac{\pi}{2}+\pi n,\qquad n\in\Z.
\]

Делим на \(10\):
\[
    \boxed{x=\frac{\pi}{20}+\frac{\pi n}{10},\qquad n\in\Z.}
\]

\section*{10. Пример с формулой \(\sin 2\alpha\)}

Решим уравнение:
\[
    8\sin 2x\cos 2x=4.
\]

Представим коэффициент \(8\) как \(4\cdot2\):
\[
    4\cdot 2\sin 2x\cos 2x=4.
\]

По формуле
\[
    2\sin\alpha\cos\alpha=\sin 2\alpha
\]
при \(\alpha=2x\):
\[
    2\sin 2x\cos 2x=\sin 4x.
\]

Получаем:
\[
    4\sin 4x=4.
\]

Делим на \(4\):
\[
    \sin 4x=1.
\]

Тогда:
\[
    4x=\frac{\pi}{2}+2\pi n,\qquad n\in\Z.
\]

Ответ:
\[
    \boxed{x=\frac{\pi}{8}+\frac{\pi n}{2},\qquad n\in\Z.}
\]

\section*{11. Формулы понижения степени}

Формулы понижения степени позволяют убрать квадрат у синуса или косинуса.

\[
    \cos^2\alpha=\frac{1+\cos 2\alpha}{2}.
\]

\[
    \sin^2\alpha=\frac{1-\cos 2\alpha}{2}.
\]

\topic{Важно:} угол увеличивается в \(2\) раза, но степень исчезает.

\section*{12. Пример с понижением степени}

Решим уравнение:
\[
    1-2\sin^2 2x=0.
\]

Используем формулу:
\[
    \sin^2\alpha=\frac{1-\cos 2\alpha}{2}.
\]

При \(\alpha=2x\):
\[
    \sin^2 2x=\frac{1-\cos 4x}{2}.
\]

Подставляем:
\[
    1-2\cdot \frac{1-\cos 4x}{2}=0.
\]

Сокращаем:
\[
    1-(1-\cos 4x)=0.
\]

Раскрываем скобки:
\[
    \cos 4x=0.
\]

Тогда:
\[
    4x=\frac{\pi}{2}+\pi n,\qquad n\in\Z.
\]

Ответ:
\[
    \boxed{x=\frac{\pi}{8}+\frac{\pi n}{4},\qquad n\in\Z.}
\]

\section*{13. Алгоритм решения}

\begin{enumerate}
    \item \important{Определите тип уравнения.} Есть ли базовый случай, формула сложения, двойной угол или квадрат?
    \item \important{Сделайте углы одинаковыми или проще.} В тригонометрии часто нужно привести выражение к одному углу.
    \item \important{Примените подходящую формулу.} Не пропускайте скобки при подстановке.
    \item \important{Сведите к базовому уравнению.}
    \[
        \sin u=a \quad \text{или} \quad \cos u=a.
    \]
    \item \important{Решите базовое уравнение.} Не забудьте параметр:
    \[
        n\in\Z.
    \]
    \item \important{Разделите обе части на коэффициент при \(x\).} Например, из
    \[
        4x=\frac{\pi}{2}+2\pi n
    \]
    получаем
    \[
        x=\frac{\pi}{8}+\frac{\pi n}{2}.
    \]
\end{enumerate}

\section*{14. Типичные ошибки}

\begin{enumerate}
    \item \important{Забыть \(n\in\Z\).} Это обязательная часть общего ответа.
    \item \important{Писать градусы вместо радиан.} В ЕГЭ по тригонометрии чаще используют радианы.
    \item \important{Потерять вторую точку на окружности.} Особенно в уравнениях \(\sin x=0\) и \(\cos x=0\).
    \item \important{Ошибиться в формуле косинуса суммы.}
    \[
        \cos(\alpha+\beta)=\cos\alpha\cos\beta-\sin\alpha\sin\beta.
    \]
    \item \important{Не поставить скобки при подстановке.} Например:
    \[
        1-2\sin^2 2x
        =
        1-2\left(\frac{1-\cos4x}{2}\right).
    \]
    \item \important{Неправильно делить общий ответ.} Если
    \[
        10x=\frac{\pi}{2}+\pi n,
    \]
    то
    \[
        x=\frac{\pi}{20}+\frac{\pi n}{10}.
    \]
    \item \important{Пытаться складывать неподобные слагаемые.} Например, \(\frac{\pi}{12}\) и \(\frac{\pi n}{6}\) нельзя объединить в одно число без учёта параметра \(n\).
\end{enumerate}

\section*{15. Краткая памятка}

\renewcommand{\arraystretch}{1.35}
\begin{table}[h!]
\caption{Формулы и базовые решения}
\centering
\begin{tabularx}{\linewidth}{@{}lX@{}}
\toprule
\textbf{Ситуация} & \textbf{Что использовать} \\
\midrule
\(\sin x=1\) & \(x=\dfrac{\pi}{2}+2\pi n,\ n\in\Z\) \\
\(\sin x=-1\) & \(x=-\dfrac{\pi}{2}+2\pi n,\ n\in\Z\) \\
\(\sin x=0\) & \(x=\pi n,\ n\in\Z\) \\
\(\cos x=1\) & \(x=2\pi n,\ n\in\Z\) \\
\(\cos x=-1\) & \(x=\pi+2\pi n,\ n\in\Z\) \\
\(\cos x=0\) & \(x=\dfrac{\pi}{2}+\pi n,\ n\in\Z\) \\
\(\cos^2\alpha-\sin^2\alpha\) & \(\cos 2\alpha\) \\
\(2\sin\alpha\cos\alpha\) & \(\sin 2\alpha\) \\
\(\sin^2\alpha\) & \(\dfrac{1-\cos2\alpha}{2}\) \\
\(\cos^2\alpha\) & \(\dfrac{1+\cos2\alpha}{2}\) \\
\bottomrule
\end{tabularx}
\end{table}

\newpage

\section*{Тренировочный блок}

\textbf{Выполните упражнения без решений.} Задачи расположены по нарастающей сложности.

\subsection*{Средний уровень}

\begin{enumerate}
    \item Решите уравнение:
    \[
        \sin x=-1.
    \]

    \item Решите уравнение:
    \[
        \cos x=0.
    \]

    \item Решите уравнение:
    \[
        1-\cos 5x\cos 2x=\sin 5x\sin 2x.
    \]
\end{enumerate}

\subsection*{Уровень выше среднего}

\begin{enumerate}[resume]
    \item Решите уравнение:
    \[
        2\sin\left(\frac{\pi}{4}-x\right)+\sqrt2\sin x=-\sqrt2.
    \]

    \item Решите уравнение:
    \[
        \cos^4 3x-\sin^4 3x=0.
    \]

    \item Решите уравнение:
    \[
        6\sin 5x\cos 5x=3.
    \]

    \item Решите уравнение:
    \[
        1-2\sin^2 3x=0.
    \]

    \item Решите уравнение:
    \[
        2\cos^2 2x-1=0.
    \]

    \item Решите уравнение:
    \[
        \cos 4x\cos x+\sin 4x\sin x=1.
    \]
\end{enumerate}

\subsection*{Сложный уровень}

\begin{enumerate}[resume]
    \item Решите уравнение:
    \[
        4\sin\left(\frac{\pi}{6}+2x\right)
        =
        2\cos 2x+2\sqrt3\sin 2x.
    \]
\end{enumerate}


\end{document}